\section{Recovering Diagnostics in a Frozen Geometry}
% -----------------------------------------------------------------------------
\subsection{Scalar linear rank-one case}
\label{subsec:rankone-recovery}
Consider again an observation $y_i=h_i^\top\theta+\varepsilon_i$ with variance $r_i$, added after a context with geometry $G_S\succ0$ and estimator $\widehat\theta_S$. Define
\begin{equation}
s_i=h_i^\top P_Sh_i,
\qquad
q_i=\frac{s_i}{r_i},
\qquad
\nu_i=y_i-h_i^\top\widehat\theta_S,
\qquad
z_i=\frac{\nu_i}{\sqrt{r_i+s_i}},
\label{eq:q-z}
\end{equation}
with $P_S=G_S^{-1}$. Then
\[
\Gamma_i=\frac{h_ih_i^\top}{r_i},
\qquad
\alpha_i=-\frac{h_i\nu_i}{r_i}.
\]
In whitened coordinates,
\[
w=G_S^{-1/2}h_i/\sqrt{r_i},
\qquad
B=ww^\top,
\qquad
b=-\frac{\nu_i}{\sqrt{r_i}}w.
\]
The unique nonzero eigenvalue of $B$ is $q_i=\norm{w}^2$, and $b$ is collinear with its eigenspace.

\begin{theorem}[Exact rank-one recovery]
In the setting above:
\begin{align}
\widehat\theta_{S+i}-\widehat\theta_S
&=\frac{P_Sh_i}{r_i+s_i}\nu_i,\label{eq:rankone-d}\\
\rho_i
&=\frac{q_i}{1+q_i}=\frac{s_i}{r_i+s_i},\label{eq:rho}\\
\mathcal I_{i\mid S}
&=\frac12\log(1+q_i)=-\frac12\log(1-\rho_i),\label{eq:rankone-info}\\
\norm{d_{i\mid S}}_{G_S}^2
&=\rho_i z_i^2,\label{eq:rankone-contextnorm}\\
\norm{d_{i\mid S}}_{G_S+\Gamma_i}^2
&=\frac{\rho_i}{1-\rho_i}z_i^2.\label{eq:rankone-postnorm}
\end{align}
The post-insertion residual is $(1-\rho_i)\nu_i$.
\end{theorem}

\begin{proof}
The first identity is Sherman--Morrison. The others follow from the single eigenvalue $q_i$ of $B$ and formulas \eqref{eq:canonical-register}--\eqref{eq:postnormd}. Details are given in Appendix~\ref{app:rankone}.
\end{proof}

\subsection{Important limit: equivalence to leverage and innovation magnitude}
When $q_i>0$, the rank-one orbit signature is entirely determined by
\begin{equation}
q_i
\qquad\text{and}\qquad
\beta_{q_i}=\norm{b}^2=q_i(1+q_i)z_i^2.
\end{equation}
For diagnostics invariant under the orthogonal action, it is therefore equivalent to the pair
\begin{equation}
\boxed{(\rho_i,|z_i|).}
\label{eq:rankone-complete-pair}
\end{equation}
Conversely,
\[
q_i=\frac{\rho_i}{1-\rho_i},
\qquad
|z_i|=\norm{b}\,[q_i(1+q_i)]^{-1/2}.
\]
The sign of $z_i$ is not contained in the orbit of $(B,b)$: in one dimension, $b$ and $-b$ are orthogonally equivalent. It must be retained separately when its interpretation matters.

If $q_i=0$, then $h_i=0$ in the identifiable geometry, so $B=b=0$ regardless of the output innovation. The parametric register is zero even though $z_i=\nu_i/\sqrt{r_i}$ may be arbitrary. This case illustrates why the output diagnostic described in the preceding section must be retained.

\begin{ochrebox}
In scalar linear regression, the invariant register does not structurally generalize the classical leverage--studentized-residual-magnitude pair: it reorganizes it and links it to the tensorial objects. Strictly non-scalar richness appears when the observation has rank greater than one, when multiple outputs or targets are considered, or when trajectories are retained.
\end{ochrebox}

\subsection{Classical diagnostics as projections}
Under the convention that residual variance is not re-estimated after deletion:
\begin{itemize}
  \item the regularized leverage of the point after insertion is $\ell_i=\rho_i$;
  \item the information gain is \eqref{eq:rankone-info};
  \item the magnitude of the studentized residual is $|z_i|$ in the orbit signature; its sign is additional output information;
  \item Cook's distance is, up to normalization by $p$, the norm in \eqref{eq:rankone-postnorm};
  \item the exact parametric influence function is $-H_S^{-1}\alpha_i$; when $H_S=G_S$, it coincides with $v_{i\mid S}^{\rm obs}$.
\end{itemize}
These relations recover the diagnostics of \citet{cook1977,hampel1974} and their local geometric interpretation \citep{cook1986,stlaurent1993}.

\begin{table}[H]
\centering
\caption{Classical diagnostic and the corresponding operation on the register.}
\small
\begin{tabularx}{\textwidth}{@{}p{0.27\textwidth}X@{}}
\toprule
Object & Projection or extension of the register \\
\midrule
Leverage & eigenvalue or trace of $C$ \\
Information gain & $\tfrac12\log\det(I+B)$ \\
Studentized-residual magnitude & norm of the innovation in the predictive metric; sign retained separately \\
Cook's distance & metric norm of the displacement $d$ \\
Parametric influence & vector $-H^{-1}\alpha$; GN approximation $-G^{-1}\alpha$ \\
Influence on a target & pairing of $v$ or $d$ with $\dd\tau$ \\
Fisher kernel & pairing of cotangent scores in a Fisher cometric \\
Experimental design & selection or aggregation of the tensors $\Gamma_i$ \\
GM-estimation & clipping or weighting policy for $(\Gamma_i,\alpha_i)$ \\
\bottomrule
\end{tabularx}
\end{table}

\subsection{Semivalues under a fixed additive geometry}
The following identities require an explicit frozen-geometry assumption.

\begin{proposition}[Static contextual factorization]
Let the tangent space be fixed, let $g_0\succ0$ be a reference metric, and let $\Gamma_j\succeq0$ be context-independent tensors. Set
\[
G_S=g_0+\sum_{j\in S}\Gamma_j.
\]
Then the marginal contribution to the intrinsic utility
\[
U_{\rm info}(S)=\frac12\log\det(g_0^{-1}G_S)
\]
is
\begin{equation}
U_{\rm info}(S+i)-U_{\rm info}(S)
=
\frac12\log\det(I+P_S\Gamma_i)
=
-\frac12\log\det(I-C_{i\mid S}).
\label{eq:static-info-factor}
\end{equation}
For any distribution $\omega_i$ over contexts, the information semivalue is the average of this contraction:
\begin{equation}
\phi_i^{\rm info,\omega}
=
\E_{S\sim\omega_i}\left[\frac12\log\det(I+P_S\Gamma_i)\right].
\label{eq:semivalue-info}
\end{equation}
\end{proposition}

Shapley value, Beta--Shapley, and Banzhaf value differ through the distribution of contexts and the chosen utility \citep{ghorbani2019,kwon2022,wang2023banzhaf}. Equation \eqref{eq:semivalue-info} shows that information Shapley is a contextual average of a scalarization of the register. Volume-based and leverage-based methods lie in a directly neighboring conceptual space \citep{xu2021volume,mendozasmith2025}.

For a predictive utility $U(S)=\tau(\widehat\theta_S)$ in a quadratic model, the marginal contribution is simply
\[
\tau(\widehat\theta_S+d_{i\mid S})-\tau(\widehat\theta_S).
\]
This identity organizes the factorization but is not, by itself, a deep theorem. In a fully retrained nonlinear model, $G_{S+i}$ is generally not $G_S+\Gamma_i$; the next section replaces the static identity by a path.

\subsection{Strictly non-scalar generalization}
When $\rank\Gamma_i>1$, scalar diagnostics no longer determine the orbit signature. Several scalarizations can be matched simultaneously.

\begin{proposition}[Collision of five scalarizations]
In $\R^3$ with $G=I$, there exist two local registers $A$ and $B$ such that
\begin{equation}
\tr C_A=\tr C_B,
\qquad
\log\det(I-C_A)=\log\det(I-C_B),
\end{equation}
\begin{equation}
\norm{d_A}=\norm{d_B},
\qquad
\norm{d_A}_{(I-C_A)^{-1}}=\norm{d_B}_{(I-C_B)^{-1}},
\qquad
t_A^\top d_A=t_B^\top d_B,
\end{equation}
but $\Sigreg_A\neq\Sigreg_B$.
\end{proposition}

\begin{proof}
Take the exact contraction spectra
\[
c_A=\left(\frac45,\frac35,\frac15\right),
\]
\[
c_B=\left(
\frac{29}{40}+\frac{3\sqrt{17}}{680},
\frac{29}{40}-\frac{3\sqrt{17}}{680},
\frac3{20}
\right).
\]
The first two components of $c_B$ are the roots of
\[
x^2-\frac{29}{20}x+\frac{893}{1700}=0.
\]
The two spectra have sum $8/5$ and satisfy $\prod_j(1-c_j)=8/125$. Choose the positive components
\[
d_A=
\left(
\frac1{10},
\sqrt{\frac{197}{500}},
\sqrt{\frac{149}{250}}
\right)
\]
and
\[
d_B=
\left(
\frac1{10},
\sqrt{
\frac{54258247}{229990400}
+
\frac{6148203\sqrt{17}}{1149952000}
},
\sqrt{
\frac{173432249}{229990400}
-
\frac{6148203\sqrt{17}}{1149952000}
}
\right).
\]
The radicands are strictly positive. Exact calculation gives
\[
d_{A,2}^2+d_{A,3}^2
=d_{B,2}^2+d_{B,3}^2
=\frac{99}{100},
\]
and
\[
\sum_j\frac{d_{A,j}^2}{1-c_{A,j}}
=
\sum_j\frac{d_{B,j}^2}{1-c_{B,j}}
=\frac{89}{50}.
\]
The decimal values used in the numerical check are therefore evaluations of these algebraic expressions, not the definition of the collision. With $t_A=t_B=e_1$, the Euclidean norms are $1$ and the target effects are $1/10$. The spectra differ, hence the signatures differ. Under a rotation $Q$, we jointly transform $C\mapsto QCQ^\top$, $d\mapsto Qd$, and $t\mapsto Qt$; the effect $t^\top d$ is therefore exactly preserved.
\end{proof}

This construction proves simultaneous non-injectivity, not merely the existence of two directions with the same norm. It is verified numerically in Section~\ref{subsec:experiments-nonscalar}.

% -----------------------------------------------------------------------------
